% from ega2-1-fr.tex
\oldpage[II]{42}
un homomorphisme $S'_{(f')}\to S_{(f)}$ que nous désignerons par
$\varphi_{(f)}$~; il lui correspond
(\hyperref[I.1.6.1-fr]{I, 1.6.1}) un morphisme
$({}^a\varphi_{(f)},\widetilde{\varphi}_{(f)}):
\operatorname{Spec}(S_{(f)})\to\operatorname{Spec}(S'_{(f')})$ de schémas
affines. Si on identifie canoniquement $\operatorname{Spec}(S_{(f)})$ au
schéma induit par $\operatorname{Proj}(S)$ sur $D_+(f)$
(\hyperref[II.2.3.6-fr]{2.3.6}), on a défini un morphisme
$\Phi_f:D_+(f)\to D_+(f')$ et ${}^a\varphi_{(f)}$ s'identifie à la
restriction de ${}^a\varphi$ à $D_+(f)$. D'autre part, il est immédiat que si
$g'$ est un second élément homogène de $S'_+$ et $g=\varphi(g')$, le
diagramme
\[
  \xymatrix{
    D_+(f)\ar[r]^{\Phi_f} & D_+(f')\\
    D_+(fg)\ar[u]\ar[r]_{\Phi_{fg}} & D_+(f'g')\ar[u]
  }
\]
est commutatif, en raison de la commutativité du diagramme
\[
  \xymatrix{
    S'_{(f')}\ar[r]^{\varphi_{(f)}}\ar[d]_{\omega_{f'g',f'}}
    & S_{(f)}\ar[d]^{\omega_{fg,f}}\\
    S'_{(f'g')}\ar[r]_{\varphi_{(fg)}} & S_{(fg)}
  }
\]
Compte tenu de la définition de $G(\varphi)$ et de
(\hyperref[II.2.3.3.2-fr]{2.3.3.2}), on voit donc que~:
\end{env}

\begin{proposition}[2.8.2]
\label{II.2.8.2-fr}
Étant donné un homomorphisme d'anneaux gradués $\varphi:S'\to S$, il existe
un morphisme et un seul $({}^a\varphi,\widetilde{\varphi})$ du préschéma
induit $G(\varphi)$ dans $\operatorname{Proj}(S')$ (dit \emph{associé à
$\varphi$} et noté $\operatorname{Proj}(\varphi)$), tel que pour tout élément
homogène $f'\in S'_+$, la restriction de ce morphisme à
$D_+(\varphi(f'))$ coïncide avec le morphisme associé à l'homomorphisme
$S'_{(f')}\to S_{(\varphi(f'))}$ correspondant à $\varphi$.
\end{proposition}

\begin{proof}
On notera qu'avec les notations précédentes, si $f'\in S'_d$, le diagramme
\[
\label{II.2.8.2.1-fr}
  \xymatrix{
    S'_{(f')}\ar[r]^{\varphi_{(f)}}\ar[d]_{\sim}
    & S_{(f)}\ar[d]^{\sim}\\
    {S'}^{(d)}/(f'-1){S'}^{(d)}\ar[r]
    & S^{(d)}/(f-1)S^{(d)}
  }
\tag{2.8.2.1}
\]
est commutatif (les flèches verticales étant les isomorphismes
(\hyperref[II.2.2.5-fr]{2.2.5})).
\end{proof}

\begin{corollary}[2.8.3]
\label{II.2.8.3-fr}
\begin{enumerate}
  \item[\rm(i)] Le morphisme $\operatorname{Proj}(\varphi)$ est affine.
  \item[\rm(ii)] Si $\operatorname{Ker}(\varphi)$ est nilpotent (et en
    particulier si $\varphi$ est injectif), le morphisme
    $\operatorname{Proj}(\varphi)$ est dominant.
\end{enumerate}
\end{corollary}

\begin{proof}
(i) est conséquence immédiate de (\hyperref[II.2.8.2-fr]{2.8.2}) et de
(\hyperref[II.2.8.1.1-fr]{2.8.1.1}). D'autre part, si
$\operatorname{Ker}(\varphi)$ est nilpotent, pour tout $f'$ homogène dans
$S'_+$, on vérifie aussitôt que $\operatorname{Ker}(\varphi_f)$ (avec
$f=\varphi(f')$) est nilpotent, donc aussi
$\operatorname{Ker}(\varphi_{(f)})$~; la conclusion résulte donc de
(\hyperref[II.2.8.2-fr]{2.8.2}) et de
(\hyperref[I.1.2.7-fr]{I, 1.2.7}).
\end{proof}

